\title{xLSTM: Extended Long Short-Term Memory}

\begin{abstract}
The foundational principles of Long Short-Term Memory (LSTM)—namely the constant error carousel and gating mechanisms—were first developed in the 1990s. LSTMs have subsequently proven to be highly durable architectures, forming the backbone of many achievements in deep learning, including the creation of early Large Language Models (LLMs). Nevertheless, the introduction of the Transformer and its efficient, parallelizable self-attention operation marked a substantial shift, surpassing the capabilities of LSTMs at large scale. In this work, we pose a straightforward inquiry: What advancements in language modeling can be realized by scaling LSTMs to billions of parameters and applying advanced techniques from contemporary LLMs while systematically addressing LSTM's established shortcomings? To this end, we first propose exponential gating augmented by dedicated normalization and stabilization strategies. Next, we alter the conventional LSTM memory design to produce: (i) sLSTM, characterized by a scalar memory, scalar update mechanism, and revised approach to memory integration, and (ii) mLSTM, a model with fully parallelizable matrix memory and a covariance-inspired update mechanism. Embedding these enhanced LSTM variants within residual block frameworks results in xLSTM blocks, which are further combined residually to assemble full xLSTM architectures. The synergistic effects of exponential gating and reengineered memory facilitate xLSTM models that compare competitively with leading Transformer and State Space Model architectures, excelling in both effectiveness and scalability.\\
Code available at: \url{https://github.com/NX-AI/xlstm}
\end{abstract}

\section{Introduction}
\label{sec:intro}

The Long Short-Term Memory (LSTM) ideas~\citep{Hochreiter:91,Hochreiter:97nips,Hochreiter:97}, i.e., the constant error carousel and gating, were introduced to overcome the vanishing gradient problem of recurrent neural networks~\citep{Hochreiter:91,Hochreiter:00book}:

\begin{align}
\label{eq:lstm_idea}
  c_t \ = \ f_t \ c_{t-1} \ + \ 
          i_t \  z_t \ , \quad  
          h_t \ = \ o_t \ \psi({c_t}) \ . 
\end{align}

The constant error carousel functions through the additive modification of the cell state $c_{t-1}$ (indicated in green) by incorporating the cell inputs $z_t$ and regulating this process using the sigmoid gates (shown in blue). The dynamics of this update are governed by the input gate $i_t$ and the forget gate $f_t$, whereas the output from the memory cell, corresponding to the hidden state $h_t$, is managed by the output gate $o_t$. Prior to producing the hidden state, the cell state undergoes normalization or is squashed via $\psi$, followed by the application of the output gate.

LSTMs have achieved significant success across multiple fields \citep{Hochreiter:01,Hochreiter:07,Schmidhuber:15} and maintained dominance in text generation applications up to the advent of Transformers in 2017~\citep{Vaswani:17}. Their utility has been validated in an array of sequential tasks, including but not limited to text generation~\citep{Graves:13,Karpathy:15blog}, handwriting synthesis~\citep{Graves:13}, sequence-to-sequence translation~\citep{Sutskever:14nips}, program analysis~\citep{Zaremba:14arxiv}, image caption creation~\citep{Karpathy:15,Hossain:19}, code generation~\citep{Karpathy:15blog}, rainfall-runoff simulations~\citep{Kratzert:18,Kratzert:19}, and hydrologic forecasting for flood warnings~\citep{Nearing:24}. Within reinforcement learning, LSTMs are recognized as top-performing sequence models, utilized by architectures such as AlphaStar for StarCraft II~\citep{Vinyals:17short}, OpenAI Five for Dota 2~\citep{OpenAI:19}, and controllers for nuclear fusion magnetic systems~\citep{Degrave:22short}. The capability of LSTMs to form abstract representations—effectively capturing and maintaining semantic knowledge in their internal states~\citep{Karpathy:15blog}—has been demonstrated, for example, through the identification of neurons specialized for number and syntax~\citep{Lakretz:19}, language properties~\citep{Bau:19}, and sentiment~\citep{Radford:17arxiv}. LSTMs continue to play an important role in contemporary, impactful settings~\citep{Degrave:22short,Nearing:24}, proving their resilience and long-term relevance.

Although LSTMs have achieved significant breakthroughs, they possess three core shortcomings: (i) They are unable to modify storage choices once made. This is demonstrated using the {\it Nearest Neighbor Search} task (additional details in Appendix~\ref{sec:appNearestNeighborSearch}), where the goal is to traverse a sequence in search of the vector most similar to a given reference, so the corresponding value can be output at the end. The left side of Figure~\ref{fig:lstmProblems} displays mean squared error results for this challenge. Traditional LSTMs cannot replace a memorized value when encountering a closer match further along the sequence; however, the proposed xLSTM addresses this issue through the use of exponential gating mechanisms. (ii) LSTMs are constrained by limited storage capacity, as they rely on scalar cell states for all stored information. This issue is illustrated with the {\it Rare Token Prediction} problem. The right side of Figure~\ref{fig:lstmProblems} presents token prediction perplexity on the Wikitext-103 dataset~\citep{Merity:17}, analyzed by token frequency. Standard LSTM architectures underperform with infrequent tokens due to these storage constraints, a challenge surmounted by our xLSTM’s use of matrix-based memory. (iii) The inherent memory interdependence between time steps—namely, the hidden-to-hidden state connections—prevents parallel computation, requiring strictly sequential operations.

These inherent limitations of LSTM have, in turn, motivated the adoption of the Transformer architecture~\citep{Vaswani:17} within language modeling. A natural question thus arises: what levels of language modeling performance can be realized when these limitations are addressed and LSTM-based models are scaled up to sizes comparable with contemporary Large Language Models?

\section{Extended Long Short-Term Memory}
\label{sec:xLSTM}

To address the shortcomings of traditional LSTM, Extended Long Short-Term Memory (xLSTM) presents two primary advancements relative to the LSTM approach described in Equation~\eqref{eq:lstm_idea}. Specifically, these advancements—exponential gating and innovative memory designs—add two new types to the LSTM family: (i) the sLSTM (refer to Section~\ref{sec:sLSTM}), which incorporates a scalar form of memory, utilizes a scalar update mechanism, and employs memory mixing; and (ii) the mLSTM (see Section~\ref{sec:mLSTM}), characterized by matrix-based memory along with a covariance-based (outer product) update procedure that supports full parallelization. Both models, sLSTM and mLSTM, leverage exponential gating to improve upon the conventional LSTM. In pursuit of efficient parallel computation, mLSTM omits memory mixing, thus removing hidden-to-hidden recurrent connections. Additionally, both variants may be scaled to accommodate multiple memory cells; in sLSTM, this scaling permits inter-cell memory mixing. The sLSTM can also be configured with multiple heads where inter-head memory mixing is not employed—mixing occurs only among cells within each head. Introducing multiple heads to the sLSTM, alongside exponential gating, establishes a distinct method for memory mixing. In contrast, for the mLSTM, assigning multiple heads or multiple cells is functionally equivalent.

When these advanced LSTM models are integrated into residual block frameworks, they form xLSTM blocks (detailed in Section~\ref{sec:xLSTMblock}). Arranging these blocks in residual layers leads to full xLSTM architectures (outlined in Section~\ref{sec:xLSTMarchitecure}). Figure~\ref{fig:xlstm_sketch} illustrates the layout of the xLSTM model and its major subcomponents.

\subsection{Review of the Long Short-Term Memory}
\label{sec:orgLSTM}

The foundational concept of the LSTM model~\citep{Hochreiter:91,Hochreiter:97nips,Hochreiter:97} involves a scalar memory cell that functions as both the primary processing and storage component, effectively mitigating the vanishing gradient issue~\citep{Hochreiter:91,Hochreiter:00book} by leveraging a constant error carousel via cell state updates. This memory cell incorporates three distinct gates: the input, output, and forget gates. The forget gate was later incorporated by \citet{Gers:00}. The following equations define the mechanism by which the LSTM memory cell updates its state at each time step $t$:

\begin{align}
c_t \ &= \  f_t \ c_{t-1} \ + \ i_t \ z_t &  & &\text{cell state} \\
h_t  \ &= \ o_t \ \tilde{h}_t \ , & \tilde{h}_t \ &= \ \psi \left( c_t\right)
  &\text{hidden state} \\
z_t \ &= \ \varphi \left( \tilde{z}_t \right) \ , 
  &\tilde{z}_t \ &=  \ \mathbf{w}^\top_{z} \ \mathbf{x}_t \ + \
  r_{z}  h_{t-1} \ + \  b_{z} \ \
  &\text{cell input} \\
i_t \ &= \ \sigma \left( \tilde{i}_t  \right) \ , 
  &\tilde{i}_t \ &= \ \mathbf{w}^\top_{i} \ \mathbf{x}_t \ + \
  r_{i}  \ h_{t-1} \ + \  b_{i} \ \
  &\text{input gate} \\
f_t \ &= \ \sigma \left(  \tilde{f}_t \right) \ , 
  &\tilde{f}_t \ &= \ \mathbf{w}^\top_{f} \ \mathbf{x}_t  \ + \
  r_{f}  \ h_{t-1} \ + \  b_{f} \ \
  &\text{forget gate} \\
o_t \ &= \ \sigma \left( \tilde{o}_t \right) \ , 
  &\tilde{o}_t  \ &= \ \mathbf{w}^\top_{o} \ \mathbf{x}_t \ + \
  r_{o}  \ h_{t-1} \ + \  b_{o} \ \
  &\text{output gate} 
\end{align}

The vectors $\mathbf{w}_{z}$, $\mathbf{w}_{i}$, $\mathbf{w}_{f}$, and $\mathbf{w}_{o}$ denote the input weights from $\mathbf{x}_t$ to the cell input, input gate, forget gate, and output gate, respectively. The terms $r_{z}$, $r_{i}$, $r_{f}$, and $r_{o}$ represent the recurrent connections from the previous hidden state $h_{t-1}$ to these same components. Bias parameters $b_{z}$, $b_{i}$, $b_{f}$, and $b_{o}$ are associated correspondingly. The activation functions $\varphi$ and $\psi$, commonly chosen as $\tanh$, are used for transforming the cell input and the hidden state. In particular, $\psi$ serves to constrain the range of the cell state, preventing it from diverging. Each gating function applies the sigmoid nonlinearity, $\sigma \left( x \right)= 1/(1+\exp(-x))$. In subsequent developments, vectors of scalar memory cells $\mathbf{c}_t \in \mathbb{R}^{d}$ were employed, which enables the utilization of recurrent weight matrices $\mathbf{R} \in \mathbb{R}^{d \times d}$, thus allowing interaction among memory cell outputs~\citep{Greff:15}; refer to Appendix~\ref{sec:apporgLSTM} for a comprehensive discussion. Investigations removing individual components confirmed that every part of the memory cell is essential for effective function~\citep{Greff:15}.

\subsection{sLSTM}
\label{sec:sLSTM}

To enable LSTMs to reconsider their memory retention choices, we incorporate exponential gates (shown in red), as well as mechanisms for normalization and stabilization. Specifically, both the input and forget gates are allowed to use exponential activation functions. For normalization purposes, we add a normalizer state that accumulates the sum of the products of the input gate and all subsequent forget gates.

The scalar sLSTM forward pass is:

\begin{align}
\label{eq:slstmforward}
c_t \ &= \  f_t \ c_{t-1} \ + \ i_t \ z_t & & 
 &\text{cell state} \\
n_t \ &= \  f_t \ n_{t-1} \ + \ i_t  & & 
 &\text{normalizer state} \\
h_t \ &= \ o_t \ \tilde{h}_t \ ,
  &\tilde{h}_t \ &= \ c_t / n_t
&\text{hidden state} \\
z_t \ &= \ \varphi \left( \tilde{z}_t \right) \ , 
  &\tilde{z}_t \ &=  \ \mathbf{w}^\top_{z} \ \mathbf{x}_t \ + \
  r_{z}  h_{t-1} \ + \  b_{z}
  &\text{cell input} \\
i_t \ &= \ \!\exp \left( \tilde{i}_t  \right) \ , 
  &\tilde{i}_t \ &= \ \mathbf{w}^\top_{i} \ \mathbf{x}_t \ + \
  r_{i}  \ h_{t-1} \ + \  b_{i}
  &\text{input gate} \\
f_t \ &= \ \sigma \left(  \tilde{f}_t \right) \ \text{OR} \
\!\exp \left(  \tilde{f}_t \right) \ ,
  &\tilde{f}_t \ &= \ \mathbf{w}^\top_{f} \ \mathbf{x}_t  \ + \
  r_{f}  \ h_{t-1} \ + \  b_{f}
  &\text{forget gate} \\
o_t \ &= \ \sigma \left( \tilde{o}_t \right) \ , 
  &\tilde{o}_t  \ &= \ \mathbf{w}^\top_{o} \ \mathbf{x}_t \ + \
  r_{o}  \ h_{t-1} \ + \  b_{o}
  &\text{output gate} 
\end{align}

We transfer the original LSTM gating techniques, i.e., input- and/or hidden-dependent gating plus bias term, to the new architectures. Exponential activation functions can lead to large values that cause overflows. Therefore, we stabilize gates with an additional state \mbox{$m_t$~\citep{Milakov:18arxiv}}:

\begin{align}
\label{eq:slstmstabil}
m_t \ &= \ \max \left( \log ( f_t ) + m_{t-1} , \log ( i_t ) \right) &\text{stabilizer state} \\
i'_t \ &= \ \!\exp \left( \log \left ( i_t \right) - m_t \right) \ = \ \!\exp \left( \tilde{i}_t - m_t \right) \ \, 
  &\text{stabil. input gate} \\
f'_t \ &= \ \!\exp \left( \log \left( f_t \right) + m_{t-1} - m_t \right) \ \, 
  &\text{stabil. forget gate}
\end{align}

As demonstrated in Appendix~\ref{sec:appsLSTM}, substituting $f_t$ with $f'_t$ and $i_t$ with $i'_t$ during the forward computation does not alter either the overall network's output or the gradients of the loss concerning the model parameters.

\paragraph{New Memory Mixing.} Like standard LSTM, sLSTM is capable of supporting several memory cells (refer to Appendix~\ref{sec:appsLSTM}).
Having multiple memory cells allows for the combination of memory content through recurrent connections $\mathbf{R}_{\mathbf{z}}$, $\mathbf{R}_{\mathbf{i}}$, $\mathbf{R}_{\mathbf{f}}$, $\mathbf{R}_{\mathbf{o}}$, which originate from the hidden state vector $\mathbf{h}$ and feed into the memory cell input $\mathbf{z}$ as well as the respective gates $\mathbf{i}$, $\mathbf{f}$, $\mathbf{o}$. 
A notable innovation in this context is the influence of exponential gating on memory mixing. The updated sLSTM architecture can contain several heads, each permitting memory mixing internally among its memory cells, but prohibiting such interaction between different heads. By integrating the concept of heads within sLSTM and incorporating exponential gating, a novel framework for memory mixing is introduced.

\subsection{mLSTM}
\label{sec:mLSTM}

To improve the storage capabilities of LSTMs, we generalize the memory cell from a single scalar $c \in \mathbb{R}$ to a matrix form $\mathbf{C} \in \mathbb{R}^{d \times d}$. Consequently, the retrieval process involves matrix multiplication. Specifically, at time $t$, we aim to encode two vectors: the key $\mathbf{k}_t \in \mathbb{R}^d$ and the value $\mathbf{v}_t \in \mathbb{R}^d$, adopting terminology from the Transformer architecture. At a subsequent time point $t + \tau$, we intend for a query vector $\mathbf{q}_{t+\tau} \in \mathbb{R}^d$ to recover the stored value $\mathbf{v}_t$. This framework aligns with the paradigm of Bidirectional Associative Memories (BAMs) \citep{Kohonen:72,Anderson:72,Nakano:72,Anderson:77}.

The covariance update rule~\citep{Sejnowski:77,Dayan:91} for storing a key-value pair is

\begin{align}
\mathbf{C}_t \ &= \ \mathbf{C}_{t-1} \ + \ \mathbf{v}_{t} \ \mathbf{k}_t^\top \ .
\end{align}

We begin by positing the application of layer normalization to the inputs prior to their projection into keys and values, ensuring these inputs possess zero mean. The process of updating the covariance matrix is regarded as optimal~\citep{Dayan:91}, as it leads to maximal separation among the retrieved binary vectors, effectively maximizing the signal-to-noise ratio. Restricting memory retrieval to interactions involving pairs and accepting the resulting quadratic computational burden—as exemplified by attention mechanisms—enables even greater separability~\citep{Krotov:16,Krotov:17,Ramsauer:21}. Furthermore, the covariance update mechanism aligns with Fast Weight Programmers~\citep{Schmidhuber:92ncfastweights,Schlag:21}, which were subsequently modified to incorporate a fixed decay rate for scaling $\mathbf{C}_{t-1}$ and a fixed learning rate for scaling 
$\mathbf{v}_{t} \mathbf{k}_t ^\top$~\citep{Ba:16}. Emulating this approach, we introduce the covariance update rule within the LSTM architecture: the forget gate is interpreted as implementing the decay rate, the input gate as representing the learning rate, and the output gate as modulating the retrieved vector.

For this matrix memory, the normalizer state is the weighted sum of key vectors, where each key vector is weighted by the input gate and all future forget gates. Again, the normalizer state keeps record of the strength of the gates. Since the dot product between query and normalizer state can be close to zero, we use the absolute value of this dot product and lower bound it by a threshold (typically 1.0) as done previously~\citep{Sun:23arxiv}. The mLSTM forward pass is:

\begin{align}
\label{eq:mlstm_recurrent_begin}
\mathbf{C}_t \ &= \  f_t \ \mathbf{C}_{t-1} \ + \ 
  i_t \ \mathbf{v}_t \ \mathbf{k}_t^\top & &  &\text{cell state} \\
\mathbf{n}_t \ &= \  f_t \ \mathbf{n}_{t-1} \ + \ 
  i_t \ \mathbf{k}_t & &  &\text{normalizer state} \\
\mathbf{h}_t  \ &= \ \mathbf{o}_t \ \odot \ \tilde{\mathbf{h}}_t
\ , \qquad \qquad 
\tilde{\mathbf{h}}_t \ = \   \mathbf{C}_t \mathbf{q}_t \ / \ 
  \max \left\{ |\mathbf{n}_t^\top \mathbf{q}_t|, 1 \right\} 
   &&&\text{hidden state} \\
\mathbf{q}_t \ &= \ \mathbf{W}_q \ \mathbf{x}_t \ + \ \mathbf{b}_q  & &  &\text{query input} \\
\mathbf{k}_t \ &= \ \frac{1}{\sqrt{d}} \mathbf{W}_k \ \mathbf{x}_t \ + \ \mathbf{b}_k  & &  &\text{key input} \\
\mathbf{v}_t \ &= \ \mathbf{W}_v \ \mathbf{x}_t \ + \ \mathbf{b}_v  & &  &\text{value input} \\
i_t \ &= \ \!\exp \!  \! \left( \tilde{i}_t  \right) \ , 
 \qquad \qquad \ \ \ \ \,
  \tilde{i}_t \ = \ \mathbf{w}^\top_{i} \ \mathbf{x}_t \ + \  b_{i} \ \ \ 
  &&&\text{input gate} \\
f_t \ &= \ \sigma \!  \left(  \tilde{f}_t \right) \ \text{OR} \ \!\exp \!  \! \left(  \tilde{f}_t \right) , \ \ \: \!
\tilde{f}_t \ = \ \mathbf{w}^\top_{f} \ \mathbf{x}_t  \ + \
  b_{f}
  &&& \text{forget gate} \\
\label{eq:mlstm_recurrent_end}
\mathbf{o}_t \ &= \ \sigma \left( \tilde{\mathbf{o}}_t \right) \ , \qquad \qquad \quad \ \ \ \,
  \tilde{\mathbf{o}}_t  \ = \ \mathbf{W}_{\mathbf{o}} \ \mathbf{x}_t \ + \
  \mathbf{b}_{\mathbf{o}}
  &&&\text{output gate} 
\end{align}

mLSTM supports several memory cells, similar to the standard LSTM architecture. Within mLSTM, employing multiple heads is functionally identical to utilizing multiple cells because no memory mixing occurs between them. To ensure stability for the exponential gating mechanism in mLSTM, we adopt the stabilization methods applied to sLSTM (refer to Equation~\ref{eq:slstmstabil}). Because memory mixing is absent in mLSTM, the recurrence relation can be restructured into a parallelizable form. Additional information can be found in Appendix~\ref{sec:appmLSTM}.

\subsection{xLSTM Architecture}
\label{sec:xLSTMarch}

\paragraph{xLSTM Blocks.}
\label{sec:xLSTMblock}
The xLSTM block is designed to generate non-linear representations of historical input data within a high-dimensional vector space, with the goal of enhancing the distinction between varied past contexts or sequences. The motivation for this separation is that, in order to accurately anticipate subsequent elements in a sequence—such as the next token—it is crucial to disentangle different historical pathways. Our architectural choices are motivated by Cover’s Theorem~\citep{Cover:65}, which posits that, by embedding data into higher dimensions through non-linear transformations, it becomes more feasible to achieve linear separability than is possible in the original feature space. To this end, we explore two forms of residual block structures: (i) a residual block featuring up-projection after non-linear operations (analogous to those used in Transformers). This version initially conducts a non-linear transformation in the base space, projects the representation linearly into a higher-dimensional domain, applies a non-linear activation, and then linearly projects it back down; refer to the left diagram in Figure~\ref{fig:Backbones}
and the third panel in Figure~\ref{fig:xlstm_sketch} for visualizations. 
An expanded illustration can be found in Figure~\ref{fig:appsLSTM_detailed}
in the appendix. (ii) Alternatively, a residual block can use up-projection prior to the non-linear step, mirroring State Space Models, in which the data first undergoes a linear transformation into a larger-dimensional space,

The model first performs a nonlinear compression of previous information within a high-dimensional space, subsequently applying a linear transformation to return data to its original dimensionality. In the context of an xLSTM block that incorporates an sLSTM, the post up-projection block is employed. Conversely, for xLSTM blocks that include an mLSTM, the pre up-projection block is selected due to the expanded memory capability in the high-dimensional domain. For a visual representation, refer to the left panel of Figure~\ref{fig:Backbones}, the third column in Figure~\ref{fig:xlstm_sketch}, or Figure~\ref{fig:appsLSTM_detailed} in the appendix.

\paragraph{xLSTM Architecture.}
\label{sec:xLSTMarchitecure}
The xLSTM architecture is devised by successively layering component blocks with residual connections, as outlined in~\citep{Srivastava:15,He:16}. This design adheres to the widely adopted pre-LayerNorm~\citep{Ba:16b} residual framework, characteristic of current Large Language Models. For additional clarification, consult the final column in Figure~\ref{fig:xlstm_sketch}.

\subsection{Memory and Speed Considerations}

In contrast to Transformers, xLSTM networks maintain linear computational costs and fixed memory usage with respect to input sequence length. The compressive nature of the xLSTM memory makes it particularly favorable for industrial scenarios and deployment on edge devices.

Unlike architectures that require parameterized memory, the memory in mLSTM operates without trainable parameters, though it incurs significant computation due to its reliance on $d \times d$ matrix memory and update operations. Here, we deliberately compromise on memory capacity to mitigate computational cost. However, the operations can be efficiently parallelized on GPUs, causing minimal delays in overall computation time.

Whereas mLSTM is amenable to parallel processing, similar to FlashAttention~\citep{Dao:22,Dao:23} and GLA~\citep{Yang:23arxiv}, the sLSTM model's memory interconnections (hidden-to-hidden) hinder parallel execution. In response, we designed a high-performance CUDA implementation, optimizing GPU memory access down to the register level. As a result, sLSTM typically achieves execution speeds within a factor of two compared to mLSTM.

\section{Related Work}
\label{sec:related}

\paragraph{Linear Attention.} Numerous approaches have been developed to address the quadratic dependency of Transformer attention on context length and to enable linear scaling with respect to sequence size. The Synthesizer model dispenses with traditional token-to-token interactions by directly learning synthetic attention patterns~\citep{Tay:20arxiv}. Linformer reduces the self-attention operation to a low-rank projection, thus allowing for a linear approximation~\citep{Wang:20arxiv}. Linear Transformer achieves a linearized version of the attention mechanism~\citep{Katharopoulos:20}. Performer employs positive orthogonal random features to efficiently approximate the softmax operation within attention using a linear method~\citep{Choromanski:21}. Alternatives to standard attention, such as fast long convolutions, are employed in Structured Global Convolution (SGConv)~\citep{Li:22} and in the Hyena Hierarchy~\citep{Poli:23}.

\paragraph{State Space Models.} In recent years, State Space Models (SSMs) have gained significant traction due to their linear scaling with sequence length and competitive results relative to Transformers. Early notable architectures in this space include the Structured State Space sequence model (S4)~\citep{Gu:21}, followed by advancements such as the Diagonal State Space (DSS) model~\citep{Gupta:22}, Gated State Space (GSS) models~\citep{Mehta:22}, S5 model~\citep{Smith:22}, Bidirectional Gated SSM (BiGS)~\citep{Wang:22}, H3 model~\citep{Fu:23}, and more recently, Mamba~\citep{Gu:24arxiv}.

\paragraph{Recurrent Neural Networks.} In recent research, Recurrent Neural Networks (RNNs) have emerged as alternative architectures to Transformer models and attention mechanisms, primarily because their computational complexity scales linearly with sequence length. Notably, RNNs utilizing Deep Linear Recurrent Units (LRUs) have demonstrated encouraging outcomes in language modeling tasks~\citep{Orvieto:23, De:24arxiv}, as have models like the Hierarchically Gated Linear RNN (HGRN)~\citep{Qin:23} and its successor, HGRN2~\citep{Qin:24arxiv}. Among the prominent RNN-based solutions for large-scale language modeling is the RWKV framework~\citep{Peng:23arxivshort,Peng:24arxivshort}, which achieves results on par with those produced by Transformer models.

\paragraph{Gating.}
Gating, a central concept in LSTM architectures, has experienced renewed attention and alternative interpretations in numerous contemporary models. Mechanisms incorporating gating are found in HGRN~\citep{Qin:23}, its successor HGRN2~\citep{Qin:24arxiv}, Gated Linear Attention (GLA)~\citep{Yang:23arxiv}, Gated State Space (GSS) architectures~\citep{Mehta:22}, Bidirectional Gated SSM (BiGS)~\citep{Wang:22}, Moving Average Equipped Gated Attention (MEGA)~\citep{Ma:22}, RWKV~\citep{Peng:23arxivshort}, as well as in Mamba~\citep{Gu:24arxiv}.

\paragraph{Covariance Update Rule.} To extend memory retention, we augmented the mLSTM cell by introducing a matrix memory governed by a covariance update rule. Similar update strategies have been leveraged by Fast Weight Programmers \citep{Schmidhuber:92ncfastweights,Schlag:21}, RWKV-5 and RWKV-6~\citep{Peng:24arxivshort}, Retention~\citep{Sun:23arxiv}, Linear Transformer~\citep{Katharopoulos:20}, and HGRN2~\citep{Qin:24arxiv}.

\paragraph{Most Related.} xLSTM is conceptually most similar to Retention~\citep{Sun:23arxiv}, RWKV~\citep{Peng:23arxivshort,Peng:24arxivshort}, and HGRN2~\citep{Qin:24arxiv}, all of which incorporate elements such as matrix memory and/or gating mechanisms. Nonetheless, unlike the recently introduced sLSTM, these existing models lack support for memory mixing. Memory mixing is vital for handling state tracking challenges, rendering LSTMs inherently more expressive than both State Space Models (SSMs) and Transformers~\citep{Merrill:24,Deletang:23}. The ability to perform state tracking is critical for tasks like executing code or following entities throughout extensive narratives.

\paragraph{Residually Stacking Architectures.} In keeping with the design principles behind most modern large-scale deep neural networks, xLSTM architectures employ the residual stacking of modular building blocks~\citep{Srivastava:15,He:16}.

This architectural approach was foundational for the development of deep convolutional neural networks~\citep{He:16} as well as the Transformer model~\citep{Vaswani:17}. Transformers have served as the central mechanism powering the rise of Large Language Models (LLMs), including systems such as GPT-3~\citep{Brown:20short}, ChatGPT~\citep{Schulman:22short}, GPT-4~\citep{Achiam:23gpt4short}, Megatron-LM~\citep{Shoeybi:19}, Gopher~\citep{Rae:21short}, ERNIE 3.0 Titan~\citep{Wang:21arxivshort}, GLaM~\citep{Du:21glamshort}, Chinese M6~\citep{Lin:21}, multilingual AlexaTM 20B~\citep{Soltan:22}, OPT~\citep{Zhang:22opt}, Chinchilla~\citep{Hoffmann:22short}, BLOOM~\citep{Scao:22arxivshort}, GLM-130B~\citep{Zeng:22arxivshort}, LaMDA~\citep{Thoppilan:22short}, PaLM~\citep{Chowdhery:22short}, Llama~\citep{Touvron:23arxiv}, and Gemini~\citep{GeminiTeam:23arxiv,Reid:24arxiv}.

\section{Experiments}
\label{sec:Exp}

In this section, we provide an empirical evaluation of xLSTM, benchmarking its performance against established approaches with an emphasis on language modeling tasks. Section~\ref{sec:ExpSyntethic} details our analysis of xLSTM's distinct capabilities using a suite of synthetic benchmarks. In Section~\ref{sec:ExpComparison}, we present a comparative study of current language modeling approaches by examining their validation perplexities, each trained on a dataset comprising 15B tokens from SlimPajama~\citep{Soboleva:23}, alongside a series of ablation experiments targeting xLSTM. Additionally, we investigate the scaling properties of all models, drawing on methodologies proposed by \citet{Kaplan:20} and \citet{Brown:20short}.

Proceeding to Section~\ref{sec:ExpLanguage}, we extend the language modeling experiments in scope. Here, xLSTM is pitted against the top-performing baselines identified in Section~\ref{sec:ExpComparison}, all of which are trained with a substantially larger data volume of 300B tokens from SlimPajama~\citep{Soboleva:23}. This assessment is organized as follows: we begin by testing the models' ability to generalize to longer context lengths, evaluate them in terms of validation perplexity and downstream performance~\citep{Sutawika:24}, and further scrutinize their effectiveness across the 571 textual domains present in the PALOMA language benchmark dataset~\citep{Magnusson:23arxivshort}. Lastly, we re-examine the scaling characteristics of each model variant, this time in a regime with twenty-fold greater training data.

\label{sec:xlstm_blocks_notation}
Throughout our experiments, the notation xLSTM[$a$:$b$] is employed to indicate an $a/b$ distribution between mLSTM-based and sLSTM-based xLSTM blocks. For instance, xLSTM[7:1] signifies a configuration of eight total blocks, comprising seven mLSTM-based blocks and a single sLSTM-based block. In the scenario where the architecture consists of 48 total blocks—a typical choice—this corresponds to 42 mLSTM-based and 6 sLSTM-based blocks. Additionally, in all settings, mLSTM and sLSTM blocks are always preceded by pre up-projection and followed by post up-projection blocks, respectively.

\subsection{Synthetic Tasks and Long Range Arena}
\label{sec:ExpSyntethic}
To begin, we evaluate the capability of xLSTM’s novel exponential gating mechanism with memory mixing on tasks involving formal languages~\citep{Deletang:23}. Subsequently, we examine how well xLSTM’s recently introduced matrix memory performs on the Multi-Query Associative Recall task~\citep{Arora:23arxiv}. Finally, we measure xLSTM’s ability to handle extended input sequences within the Long Range Arena benchmark~\citep{Tay:21}.

\paragraph{Evaluation of xLSTM's Exponential Gating Coupled with Memory Mixing.} We evaluate the proposed exponential gating with memory mixing mechanism in xLSTM, which is designed to address state tracking challenges~\citep{Merrill:24, Merrill:23}. Our study expands upon the formal language experiments introduced in \citet{Deletang:23}, allowing for training across sequences of variable lengths to support extrapolation to longer sequences. Comprehensive task descriptions and further results can be found in Appendix~\ref{sec:appExpSynthetic-formalLang}. Our benchmarking includes a range of models: Transformers, State Space Models, and various Recurrent Neural Networks. Model performance is assessed on task-specific tokens, using an accuracy metric normalized from 0 (random guessing) to 1 (exact prediction). For these tasks, we analyze 2-block variants: xLSTM[0:1] (corresponding to sLSTM alone), xLSTM[1:0] (using only mLSTM), xLSTM[1:1], alongside models such as Llama, Mamba, RWKV, Retention, Hyena, LSTM, and LSTM-in-Transformer blocks (LSTM (Block)). The experimental outcomes are displayed in Figure~\ref{fig:formal-main}. Transformer and State Space architectures lacking memory mixing—and thus lacking effective state tracking—are unable to resolve, for instance, regular grammars like the parity problem. This observation corroborates previous evidence indicating that Transformers and State Space models have inherent limitations in expressive power relative to RNNs~\citep{Merrill:24, Merrill:23, Deletang:23}.

\paragraph{Test of xLSTM's Memory Capacities on Associative Recall Tasks.} This study investigates the matrix memory mechanism in xLSTM by evaluating its memory capacity using the Multi-Query Associative Recall task~\citep{Arora:23arxiv}. In each instance, sequences consist of randomly assigned key-value pairs from a large vocabulary, all of which must be remembered for subsequent retrieval. To make the task more challenging compared to its original version, we raise both the number of key-value pairs—up to 256—and the maximum context length—up to 2048 tokens. This allows for a more comprehensive examination of memory capacities across a variety of models. Our comparison encompasses 2-block configurations for Llama, Mamba, RWKV-5, RWKV-6, xLSTM[1:1], and xLSTM[1:0]. The evaluation metric is the accuracy of recalling the correct pairs. Given that Transformers (such as Llama) possess memory capacity that grows exponentially with the model’s coding dimension~\citep{Ramsauer:21}, they serve as the benchmark for this benchmark. The comparative outcomes are displayed in Figure~\ref{fig:mqar-main}. 
Among all architectures that are not Transformer-based, xLSTM[1:1] demonstrates the highest performance, outperforming alternatives even in lower parameter regimes. Notably, the inclusion of the sLSTM block does not reduce memory capacity; rather, it appears to enhance it, particularly evident on the hardest setting with 256 key-value pairs. Further experimental details and results can be found in Appendix~\ref{sec:appExpSynthetic-mqar}, where additional extrapolation experiments suggest that the advanced memory mechanisms in xLSTM facilitate generalization to context lengths that surpass those seen during training.

\paragraph{Test of xLSTM's Long Context Capabilities on Long Range Arena.} In order to evaluate the effectiveness of xLSTM when processing extended sequences and handling substantial contextual information, we benchmark various approaches using the Long Range Arena~\citep{Tay:21}. Across all tasks, xLSTM consistently delivers robust results, indicating that its architecture is highly adept at managing the diverse challenges presented by long-context scenarios.

For more details, see Appendix~\ref{sec:appExpSynthetic-lra}.

\subsection{Method Comparison and Ablation Study}
\label{sec:ExpComparison}

In order to investigate the primary question posed in this study—namely, the performance capabilities of our novel LSTM variants when scaled for language modelling—we conduct experiments involving xLSTMs, Transformers, State Space Models, and additional baselines. Each model is trained on 15B tokens from SlimPajama within a uniform auto-regressive framework. The resulting models are evaluated on the validation set, and we further conduct ablation analyses specifically for the xLSTM architectures.

\paragraph{Comparison of xLSTM with Alternative Approaches.} To facilitate a fair comparison, all models are trained using 15B tokens sourced from SlimPajama~\citep{Soboleva:23}, with subsequent evaluation based on validation perplexity. The models under consideration include xLSTM (our proposed architecture), GPT-3 (Transformer)~\citep{Brown:20short}, Llama (Transformer)~\citep{Touvron:23arxiv}, H3 (SSM)~\citep{Fu:23}, Mamba (SSM)~\citep{Gu:24arxiv}, RWKV-4 (RNN)~\citep{Peng:23arxivshort}, RWKV-5 (RNN)~\citep{Peng:24arxivshort}, RWKV-6 (RNN)~\citep{Peng:24arxivshort}, GLA (linear Transformer)~\citep{Yang:23arxiv}, HGRN (RNN)~\citep{Qin:23}, HGRN2 (RNN)~\citep{Qin:24arxiv}, as well as RetNet (linear Transformer)~\citep{Sun:23arxiv}, Hyena (linear Transformer)~\citep{Poli:23}, and xLSTM variants (xLSTM[1:0], xLSTM[7:1]), for which definitions are provided in Section~\ref{sec:xlstm_blocks_notation}. Mixed precision was applied for model training in most cases; however, for RWKV-5, RWKV-6, GLA, and HGRN2, the mixed-precision implementation did not involve PyTorch's automated mixed precision module (details provided in Appendix Section~\ref{sec:appGenTrainProc}). The evaluated models fall into three main categories: (a) Transformers, (b) State Space Models (SSMs), and (c) Recurrent Neural Networks (RNNs), with linear Transformers included alongside RNNs due to their use of linear operations in place of conventional attention mechanisms. All model architectures are standardized to the GPT-3 configuration containing 350M parameters, specifically using an embedding dimension of 1024 and 24 residual blocks. Notably, only GPT-3 incorporates weight sharing between token and output embeddings, reducing its parameter count relative to the other models. As summarized in Table~\ref{tab:spaj15b_model_results}, xLSTM achieves lower validation perplexity than all baselines considered. Comprehensive experiment details are provided in Appendix~\ref{sec:appExpComparison}. Figure~\ref{fig:temp_sclaw15B} illustrates the scaling characteristics of this setting, suggesting that the advantage of xLSTM persists as model size increases.

\paragraph{Ablation Studies.}
As presented in Table~\ref{tab:spaj15b_model_results} and Figure~\ref{fig:temp_sclaw15B}, the xLSTM architecture demonstrates strong language modeling capabilities when trained on 15B tokens from the SlimPajama dataset. To systematically analyze the impact of modifications from the traditional LSTM to xLSTM, we progressively transition a standard LSTM model into the xLSTM configuration. The process begins by embedding LSTM layers within pre-LayerNorm residual structures. Next, this setup is expanded by incorporating a post up-projection block. In the final stage, we introduce exponential gating alongside matrix memory.

The results are shown in Table~\ref{tab:ablstudies} (top). 

Ablation experiments indicate that notable performance gains arise from incorporating both the exponential gating and the matrix memory mechanisms. Given the significant role gating plays in RNNs and State Space Models, we also evaluate various gating strategies. As shown in Table~\ref{tab:ablstudies} (bottom), our findings reveal that allowing each gate to be trainable and responsive to the input progressively enhances model effectiveness. Further details and analyses on the specific backbone elements are provided in Appendix~\ref{sec:appAblDetails}.

\subsection{xLSTM as Large Language Model}
\label{sec:ExpLanguage}

This research concludes with extensive experiments in large-scale language modeling to evaluate xLSTM’s viability as a large language model (LLM). To this end, we expand the training dataset and utilize 300B tokens from SlimPajama, an amount of data consistent with recent studies such as Mamba~\citep{Gu:24arxiv} and Griffin~\citep{De:24arxiv}.

In our comparative analysis, we benchmark xLSTM against RWKV-4, Llama, and Mamba—each exemplifying a distinct methodological category as defined in Section~\ref{sec:ExpComparison}.
RWKV-4 is chosen as the RNN baseline because suitable training precision configurations for RWKV-5, RWKV-6, and HGRN2 became available only after commencement of our 300B token runs (refer to Appendix~\ref{sec:appGenTrainProc}).
We develop models in a range of sizes (125M, 350M, 760M, 1.3B), probing all variants for their capacity to extrapolate to longer sequences and gauging performance on validation data. Additionally, we measure downstream task proficiency, assess language modeling effectiveness on 571 PALOMA benchmark domains, and systematically study scaling law phenomena.

\paragraph{Sequence Length Extrapolation.}
\label{sec:spaj300B_ctx_extrapolation}
To begin, we evaluate the ability of large-scale models (1.3B parameters) from the xLSTM, RWKV-4, Llama, and Mamba architectures to generalize to longer input sequences. The training is conducted on a maximum context length of 2048, after which the models are assessed on contexts extending up to 16384 tokens. Refer to Figure~\ref{fig:spaj300B_seq_extrapolation} for a visual summary of the outcomes. Notably, the xLSTM demonstrates consistently low perplexity scores even as the context length increases, outperforming the alternative approaches.

\paragraph{Validation Perplexity and Downstream Tasks.}
\label{sec:spaj_downstream_eval}
Additionally, across varying parameter scales, we analyze the performance of xLSTM, RWKV-4, Llama, and Mamba models on next token prediction using the SlimPajama validation dataset, alongside their effectiveness on downstream benchmarks targeting common sense reasoning abilities.

The third column in Table~\ref{tab:lmeval} shows the validation set perplexities achieved by various approaches. For all model scales, xLSTM[1:0] and xLSTM[7:1] consistently yield the lowest perplexity scores on the validation set, making them the top performers in this regard. The remaining columns in Table~\ref{tab:lmeval} summarize results on a suite of downstream tasks. xLSTM outperforms other methods on nearly all tasks and for every model size; the only notable exception is the ARC benchmark, where Mamba occasionally exhibits superior performance. Further information can be found in Appendix~\ref{sec:appExpLanguage}.

\paragraph{Performance on PALOMA Language Tasks.} Third, we evaluate the next token prediction abilities of xLSTM, RWKV-4, Llama, and Mamba across all model scales using the PALOMA language benchmarks~\citep{Magnusson:23arxivshort}. The evaluation metric employed is perplexity, assessed across 571 different text domains encompassing sources from nytimes.com to the r/depression subreddit.

Token prediction perplexity results are reported in Table~\ref{tab:ppleval}, with language modeling domains presented in the first seven columns and more detailed domain-specific benchmarks in the final five columns. On these tasks, xLSTM[1:0] demonstrates superior performance relative to xLSTM[7:1].

Specifically, xLSTM[1:0] achieves lower perplexity than Mamba in 568 out of 571 domains (99.5\%), outperforms Llama in 486 domains (85.1\%), and surpasses RWKV-4 in 570 domains (99.8\%). Additional information is provided in Appendix~\ref{sec:appExpLanguage}.

\paragraph{Scaling Laws.} Fourthly, we investigate the power-law scaling characteristics, which facilitate performance projection for increased model capacities~\citep{Kaplan:20,Brown:20short}. Figure~\ref{fig:temp_sclaw300B} illustrates these scaling trends. While all models demonstrate comparable scaling patterns, there are variations in their offsets. Among the evaluated models, RWKV-4 exhibits the poorest performance, with Llama and Mamba achieving successively better results. The xLSTM model surpasses Mamba by a margin analogous to the one separating Mamba and Llama. These observed scaling trends suggest that xLSTM is likely to maintain superior performance relative to Transformer and State-Space model families as model sizes increase.

\textbf{Generation Times and Maximal Throughput.} In our final analysis, we assess text generation latency (Figure~\ref{fig:inference_time_generation}) and peak decoding throughput (Figure~\ref{fig:inference_time_decoding_throughput}, left) across the 1.3B-parameter xLSTM models. We benchmark these results against comparably sized implementations of Mamba, Llama, and RWKV available through HuggingFace, with the Llama model employing a static key-value cache. Notably, during testing, neither fully compiled caching for the Transformer nor successful compilation of the Mamba architecture using \texttt{torch.compile} was possible. For all text generation tasks, models were evaluated using a batch size of 1 and a pre-fill of 16 tokens, a regime that particularly benefits Transformer architectures. Figure~\ref{fig:inference_time_generation} illustrates that xLSTM, along with other recurrent approaches such as Mamba and RWKV-4, exhibit linear computational complexity with respect to sequence length, in contrast to Llama’s quadratic behavior. Throughput experiments involving various batch sizes and pre-fill levels for Llama further demonstrate (see Figure~\ref{fig:inference_time_decoding_throughput}, right) that the constant memory footprint of xLSTM permits considerably larger batch processing relative to Llama, resulting in superior throughput performance.

\section{Limitations}

(i) Unlike mLSTM, the approach to memory integration in sLSTM prevents parallel processing, limiting the feasibility of highly parallelized implementations. Despite this, we have engineered a CUDA kernel optimized for sLSTM, achieving execution speeds that are currently less than twice as slow as our parallelized mLSTM baseline. (ii) The existing mLSTM CUDA kernels lack optimization, making them approximately four times slower than implementations using FlashAttention or the scan mechanism employed in Mamba. Implementing more efficient CUDA kernels, inspired by FlashAttention, should further improve speed. (iii) The matrix-based memory structure in mLSTM involves a computationally intensive $d \times d$ matrix, resulting in high computational demands. Nevertheless, parameter-free memory updating and querying can be efficiently parallelized using common matrix operations, so the added wall-clock time due to this complexity remains minimal. (iv) It is crucial to initialize the forget gates with care. (v) Because the matrix memory’s size does not depend on sequence length, longer sequences risk saturating the memory when the context window is extended. Empirically, this has not become a bottleneck for sequences up to 16k tokens (see Section~\ref{sec:spaj300B_ctx_extrapolation}). (vi) Given the significant computational expense of conducting large-scale language modeling experiments, we have not fully fine-tuned the architecture or hyperparameters, particularly for larger xLSTM models. Reaching the optimal capacity of xLSTM will likely require an intensive optimization effort.

\section{Conclusion}
We have provided a partial response to our fundamental inquiry: To what extent does language modeling improve when LSTM architectures are scaled up to billions of parameters? Our findings suggest that the performance reaches at least the level achieved by prevailing approaches such as Transformers and State Space Models. We introduced significant modifications to the LSTM, resulting in the xLSTM model through the adoption of exponential gating with memory mixing and an innovative memory mechanism. Experimental results show that xLSTM demonstrates strong performance in language modeling tasks, often surpassing leading techniques such as Transformers and State Space Models. Observed scaling laws further imply that as xLSTM models grow larger, they hold the promise of becoming formidable alternatives to contemporary Large Language Models, which are primarily constructed using Transformer-based designs. Furthermore, xLSTM's advancements could extend to fields including Reinforcement Learning, Time Series Prediction, and physical systems modeling, indicating broad applicability beyond language processing.

\end{document}